\documentclass[12pt]{article}

\usepackage[utf8]{inputenc}
\usepackage{amsmath, amssymb}
\usepackage{geometry}
\geometry{margin=1in}

\setlength{\parindent}{0pt}

\begin{document}

\begin{center}
{\Large \textbf{Solutions: Income Maintenance and Labour Supply}}\\
Chapter 3 -- Labour Market Economics
\end{center}

\vspace{1em}

\textbf{Given:} $G=800$ per week, $t=0.4$, $w=25$ per hour, total available time $T=60$ hours/week.

While eligible for benefits:
\[
C = G + (1-t)wh
\quad \text{for } wh \leq \frac{G}{t}.
\]
After benefits are exhausted:
\[
C = wh.
\]

\section*{(a) Break-even earnings and hours (3 marks)}

\subsection*{Step 1: Break-even level of earnings}

Benefits fall to zero when the amount clawed back equals the guarantee.  
The clawback on earnings $wh$ is:
\[
t(wh).
\]
Break-even occurs when:
\[
t(wh) = G.
\]
Solve for earnings:
\[
wh = \frac{G}{t} = \frac{800}{0.4} = 2000.
\]
\textbf{Answer:} Break-even earnings are \textbf{\$2000 per week}.

\subsection*{Step 2: Hours at which benefits are exhausted}

Convert break-even earnings into hours by dividing by the wage:
\[
h^{BE} = \frac{2000}{25} = 80.
\]

\textbf{Interpretation:} The model implies the worker would need to work \textbf{80 hours} for benefits to be fully exhausted.  
But the time endowment is only $T=60$ hours, so $h \le 60$. Therefore:

\begin{itemize}
\item The worker \textbf{cannot reach} the break-even point within the available time.
\item With a 60-hour time endowment, the worker remains \textbf{on benefits for all feasible hours}.
\end{itemize}

\textbf{Answer:} Benefits would be exhausted at \textbf{80 hours}, which is \textbf{not feasible} given $T=60$.

\section*{(b) Net wage and effective marginal tax rate (2 marks)}

\subsection*{Step 1: Effective net wage while on benefits}

While receiving benefits:
\[
C = G + (1-t)wh.
\]
The slope of $C$ with respect to $h$ is the effective net wage:
\[
\frac{dC}{dh} = (1-t)w.
\]
Compute:
\[
(1-t)w = (1-0.4)\cdot 25 = 0.6 \cdot 25 = 15.
\]
\textbf{Answer:} Effective net wage is \textbf{\$15 per hour} while on benefits.

\subsection*{Step 2: Effective marginal tax rate (EMTR)}

The EMTR is the fraction of an additional dollar of earnings that is lost due to benefit reduction.  
Since benefits are reduced at rate $t$, the EMTR from the program is:
\[
EMTR = t = 0.4 = 40\%.
\]
Equivalently, the worker keeps $(1-t)=0.6$ of each additional dollar earned.

\textbf{Answer:} Effective marginal tax rate is \textbf{40\%}.

\section*{(c) Budget constraint in $(h,C)$ space (3 marks)}

\subsection*{Key points and slopes}

\begin{itemize}
\item At $h=0$:
\[
C = G = 800.
\]
So the vertical intercept is $(0,800)$.

\item While on benefits, the budget line is:
\[
C = 800 + 15h,
\]
because $(1-t)w=15$.

\item The break-even point would be at:
\[
h^{BE}=80, \quad C^{BE}=wh^{BE}=25\cdot 80=2000,
\]
but \textbf{this is outside the feasible set} since $h \le 60$.

\item Maximum feasible hours: $h=T=60$.
Disposable income at $h=60$ (still on benefits) is:
\[
C(60) = 800 + 15(60) = 800 + 900 = 1700.
\]
\end{itemize}

\subsection*{How to draw the graph}

\begin{enumerate}
\item Put $h$ on the horizontal axis and $C$ on the vertical axis.
\item Plot the point $(0,800)$.
\item Draw a straight line from $(0,800)$ with slope $15$ up to $h=60$.
\item Mark the endpoint at $(60,1700)$.
\item Optionally, show the \emph{hypothetical} kink at $h=80$ where the slope would switch from $15$ to $25$, but indicate it is beyond $T=60$ and thus not attainable.
\end{enumerate}

\textbf{Slopes to label:}
\begin{itemize}
\item While on benefits: slope $=(1-t)w = 15$.
\item After benefits exhausted: slope $=w=25$ (not reached here).
\end{itemize}

\section*{(d) Extensive vs intensive margin effects (3 marks)}

\subsection*{Extensive margin (participation: work or not)}

\begin{itemize}
\item The guarantee $G=800$ raises income at $h=0$ relative to a no-program world.
\item This creates an \textbf{income effect} that makes non-work more attractive for some individuals.
\item Therefore, the program can reduce labour force participation: some individuals may choose $h=0$ when they would have worked without the guarantee.
\end{itemize}

\subsection*{Intensive margin (hours conditional on working)}

\begin{itemize}
\item While on benefits, the effective net wage is reduced from $25$ to $15$.
\item The lower net wage reduces the opportunity cost of leisure, generating a \textbf{substitution effect} toward more leisure and fewer hours.
\item The guarantee also increases non-labour resources, which (if leisure is normal) produces an \textbf{income effect} toward more leisure.
\item Both effects typically reduce hours worked for those in the benefit range.
\end{itemize}

\textbf{Important note for this parameter set:} Because $h^{BE}=80>T=60$, the worker is on benefits for all feasible hours. That means the worker faces the reduced net wage for the entire feasible range, strengthening the intensive-margin disincentive.

\section*{(e) Increase $t$ from 0.4 to 0.7 (4 marks)}

Now set $t=0.7$ holding $G=800$ and $w=25$ fixed.

\subsection*{(e1) How the budget constraint changes}

\begin{itemize}
\item New effective net wage while on benefits:
\[
(1-t)w = (1-0.7)\cdot 25 = 0.3\cdot 25 = 7.5.
\]
So the benefit-region budget line becomes much flatter:
\[
C = 800 + 7.5h.
\]

\item New break-even earnings:
\[
wh = \frac{G}{t} = \frac{800}{0.7} \approx 1142.86.
\]
New break-even hours:
\[
h^{BE} = \frac{1142.86}{25} \approx 45.71.
\]
This time, the break-even point is \textbf{feasible} because $45.71 < 60$.

\item Therefore, the budget constraint is now \textbf{kinked within the feasible set}:
\begin{itemize}
\item For $0 \le h \le 45.71$, slope $=7.5$ (on benefits).
\item For $h \ge 45.71$, benefits are exhausted and slope returns to $w=25$.
\end{itemize}
\end{itemize}

\subsection*{(e2) Predicted effect on labour supply}

\begin{itemize}
\item \textbf{Intensive margin:} Increasing $t$ lowers the net wage in the benefit range (from 15 to 7.5), which strengthens the substitution effect away from work. This tends to reduce hours among individuals who remain in the benefit range.
\item However, because the kink now occurs within feasible hours, some individuals may adjust hours to locate near the kink (bunching), depending on preferences.
\item \textbf{Extensive margin:} A higher $t$ reduces the reward from entering work at low hours (since the net wage is lower), which can reduce participation for individuals near the margin of working.
\end{itemize}

\subsection*{(e3) Equity--efficiency trade-off}

\begin{itemize}
\item \textbf{Equity:} A higher $t$ targets benefits more strongly to low earners by reducing benefits more rapidly as earnings rise. This can reduce spending on higher-earning recipients and focus support on those with the least income.
\item \textbf{Efficiency:} The cost is weaker work incentives. The higher effective marginal tax rate discourages additional work in the phase-out region and may reduce employment and hours.
\item Policymakers must balance the goal of income support (insurance and redistribution) against the distortions from high EMTRs.
\end{itemize}

\vfill
\textbf{End of Solutions}

\end{document}
